% main.tex — The Limitations of Generative AI (working draft)
% Build with `make` (requires TeX Live with latexmk and biber).
\documentclass[11pt,a4paper]{article}

\usepackage[margin=2.5cm]{geometry}
\usepackage{amsmath,amssymb}
\usepackage{xcolor}
\usepackage{microtype}
\usepackage{booktabs}
\usepackage[style=authoryear,backend=biber,maxcitenames=2]{biblatex}
\usepackage[colorlinks=true,linkcolor=blue!50!black,citecolor=blue!50!black,urlcolor=blue!50!black]{hyperref}

\addbibresource{references.bib}

% Marker for passages that still need substantive work.
\newcommand{\draft}[1]{\textcolor{red}{[DRAFT: #1]}}

\title{The Limitations of Generative AI\\[0.3em]\large A critical working draft}
\author{Nadia Yvette}
\date{\today}

\begin{document}
\maketitle

\begin{abstract}
This write-up systematizes the known limitations of contemporary generative AI
systems---large language models chief among them---and distinguishes carefully
between what is \emph{observed}, what is \emph{explained}, and what is merely
\emph{narrated}. The organizing claim is that the most consequential failures
are not incidental bugs but structural consequences of how these systems are
built: next-token prediction over compressed corpora, optimization for
plausibility rather than truth, and evaluation regimes that reward fluent
imitation of understanding.
\draft{Sharpen the thesis once the section drafts stabilize.}
\end{abstract}

\tableofcontents

\section{Introduction}
\label{sec:introduction}
\draft{Motivate the write-up; state scope, method, and intended audience.}

Generative systems now produce fluent text, images, audio, and code at near-zero
marginal cost. Fluency, however, is not knowledge, and plausibility is not
correctness \autocite{bender2021stochastic,bender2020climbing}. This document
catalogs the failure modes that follow from that gap, with attention to which
limitations are \emph{fundamental} to the paradigm and which are
\emph{engineering shortfalls} that may yield to better tooling, data, or
verification.

\section{Background and terminology}
\label{sec:background}
\draft{Define the object of study: autoregressive LLMs, diffusion models,
multimodal generators. Fix vocabulary early (hallucination vs.\ confabulation,
grounding, calibration).}

The dominant architecture for language generation is the transformer
\autocite{vaswani2017attention}, trained to continue sequences. Image and audio
generators increasingly share the denoising/diffusion paradigm. Throughout, we
use \emph{limitation} to mean a systematic constraint on reliability,
grounding, or reasoning---not a one-off defect of a particular product.

\section{A taxonomy of limitations}
\label{sec:taxonomy}
\draft{This is the core of the write-up; each subsection gets its own
evidence base before any generalizing claims.}

\subsection{Reliability: confabulation and hallucination}
\label{sec:hallucination}
Generative models produce statements unsupported by their inputs, their
training data, or the world
\autocite{ji2023survey,mallen2023trust}. The phenomenon is not rare-model
noise: it follows from training objectives that reward plausible continuations.

\subsection{Reasoning, compositionality, and planning}
\label{sec:reasoning}
Transformers degrade sharply on tasks whose solution requires reliable
composition of sub-computations \autocite{dziri2023faith}, and large language
models demonstrably fail classical planning benchmarks
\autocite{valmeekam2023planning,kambhampati2024llms}.

\subsection{Grounding: form versus meaning}
\label{sec:grounding}
Language-only training leaves meaning underdetermined: the same strings are
compatible with many worlds \autocite{bender2020climbing}. Multimodal training
narrow but does not close this gap.

\subsection{Data-derived limits}
\label{sec:data}
Bias, temporal cutoffs, duplication, contamination of evaluation sets, and the
thin coverage of the long tail propagate directly into model behavior
\autocite{bommasani2021opportunities,weidinger2021ethical}.

\subsection{Calibration and uncertainty}
\label{sec:calibration}
Models often lack well-calibrated confidence over their own outputs; fluent
delivery is a poor proxy for posterior probability
\autocite{kadavath2022language}.

\subsection{Evaluation pathologies}
\label{sec:evaluation}
Heuristic shortcuts \autocite{mccoy2019right}, benchmark contamination, and
Goodhart effects make headline metrics unreliable guides to capability.

\subsection{Social and systemic risks}
\label{sec:social}
Harms to persons and institutions \autocite{weidinger2021ethical}, and the
compounding of the above when generators are chained into agentic pipelines.

\subsection{Memory and context: the read side of state}
\label{sec:memory}
\draft{Expand from notes/memory-and-diffusion.md; keep the read/write split
explicit.}

Contemporary models do not remember: they are functions of their context
window, and long-term memory is externalized to engineering scaffolding.
The \emph{write} side --- changing what the system knows through
experience --- is the dual-control problem of
Section~\ref{sec:learning-in-the-loop}. The \emph{read} side is a distinct
failure that persists even with frozen weights: models reliably use
information at the edges of long contexts but miss the middle
\autocite{liu2024lost}, and models advertising 32K-plus token windows often
fail realistic tasks well below their advertised limits
\autocite{hsieh2024ruler}. Context is best understood as short-term working
memory with unmeasured fidelity --- and retrieval scaffolds inherit the same
read-reliability problem one layer out.

\subsection{Learning in the loop: dynamical failures}
\label{sec:learning-in-the-loop}
\draft{Expand from notes/online-neural-control.md; keep the dynamical-versus-statistical distinction sharp.}

The preceding failures are \emph{statistical}: errors average out over
samples. A distinct family arises when a network is updated online inside
its own control loop, where errors \emph{circulate through feedback}.
Combining bootstrapped targets with function approximation and off-policy
data diverges --- provably, even with linear approximation and exact
dynamic-programming updates \autocite{baird1995residual,
tsitsiklis1997analysis,vanhasselt2018deadly}. The acting controller
additionally corrupts its own data distribution: agents overfit their
earliest experiences \autocite{nikishin2022primacy}, shifted distributions
overwrite prior learning \autocite{french1999catastrophic}, and prolonged
continual training destroys the \emph{capacity to learn} altogether
\autocite{dohare2024plasticity}. A learned dynamics model invites
optimizing against its own flattering extrapolations
\autocite{janner2019trust}. The classical umbrella is Feldbaum's dual
control, intractable even for linear-Gaussian systems
\autocite{wittenmark1995adaptive}; the principled escape is two-timescale
separation of estimation from action \autocite{borkar1997stochastic}.

\section{Alternatives in the generator design space}
\label{sec:alternatives}
\draft{Expand from notes/memory-and-diffusion.md; keep the plain-language
framing for the intended audience.}

Autoregressive models write like a typist, one token after another;
diffusion models sculpt the whole answer at once, refining noise into text
like a photograph developing. The sculptor is measurably faster --- over a
thousand tokens per second on current hardware, with sub-second first
tokens \autocite{inception2025mercury} --- and handles sound and images
natively; an 8B-parameter diffusion model trained on a fraction of the data
matches an equivalent autoregressive model \autocite{nie2025llada}. The
typist still holds two advantages: very long documents are cheaper to serve
thanks to key--value caching, and deliberate step-by-step reasoning is
currently built on sequential generation. Better on speed and modality;
better overall --- still an open, empirical question.

\section{What the limitations imply}
\label{sec:implications}
\draft{Survey the mitigation landscape---retrieval augmentation, tool use,
verification, neuro-symbolic hybrids---and assess honestly which limitations
each addresses and which it merely hides \autocite{gao2023rag}.}

\subsection{The neurosymbolic case: governance demands structure}
\label{sec:neurosymbolic}
\draft{Expand from notes/neurosymbolic-governance.md; keep the coherence
criterion and the corrected deployment claim.}

Some obligations on AI systems are not statistical questions. Privacy, in
its dominant positive theory (contextual integrity), is a system of norms
over who may transmit what information about whom, to whom, in which
context --- and when it was formalized, that theory became a deontic logic
with temporal operators \autocite{barth2006privacy}. Reading the formal
requirements off yields exactly multi-type, multi-agent, multimodal logic
with temporality and probabilism for risk: recipient beliefs, role and
group structure, retention over time, leak-path probabilities. A corpus
gradient can imitate descriptions of norms; it cannot apply one. It can
produce instances of protective behavior, but coherence --- norm
consistency, auditability, reversibility --- requires inspectable
structure. The semiotic layer supplies what the formalization leaves
informal: a typed account of what kind of sign is transmitted about whom
(indexical links, iconic portrayals, symbolic misuses). Narrow hybrids are
deployed commercially; the full normative profile is deployed nowhere
\autocite{garcez2023neurosymbolic,trinh2024alphageometry,
hubert2025alphaproof}. Architecturally, the precedent is the ethical
governor \autocite{arkin2009governing}: normative state outside the
trained weights, neural perception fast and graded --- the same
two-timescale separation that Section~\ref{sec:learning-in-the-loop}
derives from stability, derived here from accountability.

\subsection{The shape of the machine shapes the science}
\label{sec:saas}
\draft{Expand from notes/saas-political-economy.md; keep intent deliberately
underdetermined; the effects are the citable part.}

One further limitation is not inside the models but around them. The
dominant delivery mechanism of the past decade --- models offered only as
hosted services, weights withheld --- was selected under incentives where
keeping the service proprietary had commercial value, and the choice has
measurable costs for research: transparency scores for major developers
are low and declining, with the worst indicators exactly where a service
model hides things (training data, compute)
\autocite{bommasani2023fmti}; published findings silently break when
hosted models are updated or retired without documentation
\autocite{pozzobon2023blackbox}; and access to the necessary scale was
already concentrating industrially before the commercial boom
\autocite{ahmed2020dedemocratization,kapoor2023leakage}. An internal
document leaked from Google in 2023 acknowledged the moat reasoning
directly \autocite{patel2023nomoat}. Competing explanations --- safety,
legal exposure --- predict much the same behavior, so motive must remain
an open question; the measurable damage to the research commons is the
part that can be stated without hedging.

\section{Open questions}
\label{sec:open-questions}
\draft{Collect open research questions, including which limitations are
in-principle vs.\ in-practice.}

\section{Conclusion}
\label{sec:conclusion}
\draft{Summarize; avoid the twin failure modes of hype and dismissiveness.}

\printbibliography

\end{document}
